
\chapter{Trabalhos Relacionados}
\label{cap:trabalhos-relacionados}

Os trabalhos mais próximos da proposta geram modelos BPMN a partir de texto com modelos de linguagem e medem pelo menos uma das grandezas usadas aqui: a proporção de saídas válidas, a redução de \textit{tokens} obtida por uma representação intermediária ou a qualidade do modelo gerado. Eles estão agrupados em três linhas: a geração direta do XML, a geração por meio de uma representação intermediária e o ajuste de modelos pequenos para a tarefa. A Seção~\ref{sec:tr-comparacao} reúne os números de cada trabalho e os da proposta na Tabela~\ref{tab:tr-comparacao} e discute até onde eles podem ser comparados.

\section{Geração Direta de BPMN XML}
\label{sec:tr-xml}

\citeonline{lauer2026assessing} avaliaram 17 modelos de linguagem abertos, de 1 a 235 bilhões de parâmetros, na geração de BPMN XML a partir de 105 pares de texto e modelo. Cada modelo recebeu um \textit{prompt} com um exemplo de saída, temperatura de 0,1 e cinco execuções por texto. Quando o XML gerado era inválido, o modelo recebia os erros e tinha uma única chance de corrigi-lo. A validade foi medida da mesma forma que na proposta, pela validação contra o esquema XSD do BPMN 2.0. Seis modelos, quase todos com menos de 7 bilhões de parâmetros, produziram tão poucos arquivos válidos que foram excluídos da análise. Entre os 11 restantes, a validade variou de 30,7\%, no Falcon3-10B, a 97,3\%, no Llama-3.3-70B, e os modelos Qwen2.5 de 14 e de 32 bilhões de parâmetros ficaram em 54,7\% e 51,1\%. Os autores concluem que gerar BPMN XML válido continua sendo um desafio para os modelos atuais, e que, nos aspectos semânticos, os modelos de linguagem ficam abaixo de modeladores humanos.

Esses resultados servem de referência externa para a validade da geração direta. O A1, em que o DeepSeek~V4~Pro escreve o XML em uma única tentativa, teve 97,8\% de validade, no mesmo patamar do melhor modelo de \citeonline{lauer2026assessing}, que contava com uma rodada de correção. A dificuldade apontada pelos autores aparece, assim, sobretudo nos modelos menores, e é ela que o ajuste do modelo pequeno procura contornar.

\section{Geração por Representação Intermediária}
\label{sec:tr-representacoes}

Uma alternativa à geração direta é pedir ao modelo uma representação mais simples, convertida depois em BPMN por um programa. O ProMoAI, de \citeonline{kourani2024promoai}, pede a um modelo de grande porte um código Python que constrói um modelo POWL, depois traduzido para BPMN, e foi avaliado de forma qualitativa. A representação \textit{JSON branches}, de \citeonline{kopke2024efficient}, descreve os ramos do processo em JSON, com o objetivo declarado de reduzir o número de \textit{tokens} e facilitar a adesão ao esquema. A representação \textit{Process Model Elements} (PME), de \citeonline{voelter2024multimodal}, descreve o modelo como listas de elementos e fluxos.

\citeonline{brissard2025pmr} compararam nove dessas representações, incluindo o próprio BPMN XML, sobre o PMo, o mesmo conjunto usado no teste da proposta. Na primeira parte do estudo, mediram o tamanho dos modelos de referência em cada representação. Em relação ao BPMN XML completo, com as coordenadas do diagrama, a remoção das coordenadas reduz o número de \textit{tokens} em 70\%, e as representações compactas chegam a reduções de 88\%, no \textit{JSON branches}, e de 93\%, no Mermaid. Na segunda parte, usaram o Llama-3.3-70B, com temperatura de 0,2, para gerar os modelos em cada representação, restringindo os elementos a tarefas, eventos de início e fim, \textit{gateways} exclusivos e paralelos e fluxos de sequência. A qualidade foi medida pela similaridade PME, que alinha os elementos do modelo gerado aos da referência pela semelhança dos rótulos e resume a concordância por um coeficiente de Dice. A similaridade geral ficou entre 0,43, no BPMN XML, e 0,54, na representação textual em ramos, e o código POWL ficou em 0,27, em parte porque 40\% dos modelos gerados nessa representação eram inválidos. Em todas as representações, os modelos gerados tinham, em média, cerca de oito nós a menos que a referência, com perda maior nos \textit{gateways}. O estudo não reporta a validade das demais representações nem o custo da geração.

\citeonline{licardo2026bpmn} compararam diretamente as duas estratégias. No BPMN Assistant, o modelo gera um JSON hierárquico que um programa converte em BPMN 2.0 XML, e a comparação é feita com o mesmo modelo gerando o XML. Foram avaliados nove modelos, entre eles o GPT-5.1, o Claude~4.5~Sonnet, o Llama-3.3-70B e o Qwen2.5-72B, em 60 descrições de processos com sete ou oito atividades, geradas pelo GPT-4.1 e modeladas por anotadores. Uma geração conta como falha quando o modelo não pode ser desenhado, por erro de sintaxe ou por fluxo que aponta para um elemento inexistente. Houve 2 falhas em 540 gerações com o JSON e 11 com o XML. A similaridade com a referência, medida pela distância de edição em grafos relativa após a normalização dos rótulos por outro modelo de linguagem, foi de 0,72 com o JSON e de 0,70 com o XML, uma diferença que os autores consideram mínima. O ganho do JSON esteve na eficiência: 688 \textit{tokens} de saída, contra 1.832 do XML, e latência média de 13,4 segundos, contra 24,8. Em troca, a entrada passou de 474 para 2.679 \textit{tokens}, porque o \textit{prompt} precisa descrever o formato JSON.

O desenho de \citeonline{licardo2026bpmn} é o mais próximo da comparação entre o A1 e o A2: o mesmo modelo, ora gerando o XML, ora uma representação intermediária. Os dois estudos encontram a mesma troca nos \textit{tokens}, com saída menor e entrada maior, porque a descrição do formato vai no \textit{prompt}. Divergem na validade. No BPMN Assistant, a representação intermediária reduziu as falhas, e na proposta a DSL as aumentou, de 2,2\% para 37,7\% das gerações do DeepSeek~V4~Pro (Seção~\ref{sec:res-validade}). Uma explicação possível está nas próprias representações: o JSON é um formato amplamente presente nos dados de treino dos modelos, enquanto a BPMN-DSL é uma notação nova, aprendida apenas pela gramática no \textit{prompt}, e cuja primeira versão exigia declarar identificadores antes de saber se seriam usados. Os processos do BPMN Assistant também são menores e só usam \textit{gateways} exclusivos e paralelos.

\section{Modelos Pequenos Ajustados para a Tarefa}
\label{sec:tr-slm}

\citeonline{li2025llmbased} compararam cinco estratégias para gerar BPMN na linguagem DOT, do Graphviz, a partir do MaD \cite{li2023mad}, um conjunto de cerca de 30 mil pares de texto e modelo em 15 categorias de processos: \textit{prompt} sem exemplos, com exemplos, as duas versões com raciocínio passo a passo, e o GPT-3.5-turbo ajustado pela interface de ajuste do provedor. Nas 3.000 descrições de teste, a taxa de saídas com erro de sintaxe caiu de 6,1\%, sem exemplos, para 0,43\%, com exemplos e raciocínio, e para zero no modelo ajustado. O modelo ajustado também teve a maior similaridade com a referência, de 1,000 nos nós e de 0,944 no grafo, medidas por semelhança de rótulos e de representações vetoriais do grafo. \citeonline{soliman2025size} seguiram a mesma linha com um modelo aberto, ajustando o Qwen2.5-Coder-14B no MaD para gerar DOT, convertido depois em BPMN XML por um programa. O mesmo modelo base foi ajustado também para gerar o BPMN diretamente, com um conjunto de treino próprio, e as duas vias foram comparadas. A saída em DOT teve cerca de 55 linhas nos processos de complexidade média e 77 nos complexos, contra 241 e 341 do BPMN. Os autores reportam acurácia, verificada por isomorfismo entre o grafo gerado e o de referência, de 97,3\% para o DOT e de 88\% para o BPMN, mas o texto não deixa claro se os 88\% se referem à geração direta ou ao BPMN obtido pela conversão do DOT, nem sobre quantos itens do conjunto de teste as duas taxas foram calculadas.

\citeonline{celikmasat2026instrubpm} ajustaram, com LoRA de posto 16, o Qwen3-4B, de 4 bilhões de parâmetros, em cerca de 21,5 mil pares do MaD, também com saída em DOT. Nas 180 descrições de teste, retiradas do mesmo conjunto, 99,4\% das saídas do modelo ajustado puderam ser lidas pelo analisador de DOT. A semelhança estrutural com a referência, medida pela distância de edição em grafos relativa, que compara só a topologia e ignora os rótulos, foi de 99,4 em 100, contra 43,6 do melhor modelo comercial avaliado, o Gemini~2.5~Pro, e 38,5 do Qwen2.5-Coder-7B sem ajuste, o mesmo modelo base usado no A4.

Esses resultados sustentam a premissa de \citeonline{belcak2025small}, de que modelos pequenos especializados bastam para tarefas repetitivas e bem delimitadas. A mesma conclusão aparece fora da geração de processos. Na classificação de 16 mil petições iniciais e denúncias do Tribunal de Justiça do Ceará em cinco classes processuais, um BERT ajustado alcançou F1 macro de 0,88 \cite{aguiar2021text}, e a combinação com modelagem de tópicos e com as leis citadas nos documentos chegou a 0,89 \cite{aguiar2022topic}. Nos trabalhos com o MaD, porém, o teste vem do mesmo conjunto que o treino, com os mesmos domínios e o mesmo estilo de escrita. Na proposta, o treino vem quase todo do \textit{handbook} da GitLab, e o teste, do PMo, que reúne processos de outras fontes. Nessa condição, o modelo ajustado teve 87,0\% de validade, abaixo do que os trabalhos com o MaD reportam, e fidelidade inferior à dos modelos de grande porte (Seção~\ref{sec:res-modelo-especializado}).

\section{Compressão de \textit{Tokens} e Garantia de Validade}
\label{sec:tr-compressao}

A redução do número de \textit{tokens} em sistemas com modelos de linguagem tem sido tratada sobretudo do lado da entrada. O LLMLingua \cite{jiang2023llmlingua} usa um modelo auxiliar para remover do \textit{prompt} os trechos que menos contribuem para a resposta, com pouca perda de desempenho. Essa compressão aceita perder informação e não reduz a saída, que, na geração de marcação, é a parte mais longa e a mais cara (Seção~\ref{subsec:tokenizacao}). As duas abordagens são, portanto, independentes e podem ser combinadas.

A decodificação restrita \cite{poesia2022synchromesh} é a principal alternativa para garantir a validade sintática. Ela impede, durante a geração, que o modelo escreva uma sequência fora da gramática, e poderia ser aplicada tanto ao XML quanto à DSL. Ela não reduz o número de \textit{tokens} gerados, e uma gramática livre de contexto não impede que um fluxo aponte para um elemento que não existe. Essa referência a elemento não declarado foi a principal causa de invalidade nos braços de DSL do DeepSeek~V4~Pro, do GLM~5.2 e do modelo ajustado (Seção~\ref{sec:res-validade}). Nenhum dos trabalhos analisados neste capítulo a utiliza.

\section{Comparação com a Proposta}
\label{sec:tr-comparacao}

A Tabela~\ref{tab:tr-comparacao} resume os trabalhos discutidos e os principais braços da proposta.

\begin{table}[htb]
	\Caption{\label{tab:tr-comparacao} Trabalhos relacionados e proposta}%
	\IBGEtab{}{%
		\footnotesize
		\setlength{\tabcolsep}{3pt}
		\begin{tabular}{>{\raggedright\arraybackslash}p{2.5cm} >{\raggedright\arraybackslash}p{2.5cm} >{\raggedright\arraybackslash}p{1.7cm} >{\raggedright\arraybackslash}p{1.8cm} >{\raggedright\arraybackslash}p{2.2cm} >{\raggedright\arraybackslash}p{1.5cm} >{\raggedright\arraybackslash}p{2.2cm}}
			\toprule
			Trabalho & Modelo e adaptação & Saída & Teste & Validade & Redução de \textit{tokens} & Qualidade \\
			\midrule \midrule
			\citeonline{lauer2026assessing} & 11 modelos abertos, de 8B a 235B, \textit{prompt} & BPMN XML & 105 pares & 30,7\% a 97,3\% (XSD) & --- & semântica de 0,53 a 0,58 \\
			\midrule
			\citeonline{brissard2025pmr} & Llama-3.3-70B, \textit{prompt} & 9 representações & PMo, 55 & POWL: 60\% & 70\% a 93\%\textsuperscript{a} & PME de 0,27 a 0,54 \\
			\midrule
			\citeonline{licardo2026bpmn} & 9 modelos de grande porte, \textit{prompt} & JSON; BPMN XML & 60 sintéticos & 99,6\%; 98,0\% (desenhável) & 62\% & 1 $-$ RGED: 0,72; 0,70 \\
			\midrule
			\citeonline{li2025llmbased} & GPT-3.5-turbo, \textit{prompt} e ajuste & DOT & MaD, 3.000 & 93,9\% a 100\% (sintaxe) & --- & nós 1,000; grafo 0,944\textsuperscript{b} \\
			\midrule
			\citeonline{soliman2025size} & Qwen2.5-Coder-14B, ajuste & DOT; BPMN XML & MaD, subconjunto não informado & --- & cerca de 77\%\textsuperscript{c} & acurácia 97,3\%; 88\%\textsuperscript{d} \\
			\midrule
			\citeonline{celikmasat2026instrubpm} & Qwen3-4B, LoRA & DOT & MaD, 180 & 99,4\% (sintaxe) & --- & R-GED 99,4\textsuperscript{b} \\
			\midrule
			Proposta, A1 & DeepSeek~V4~Pro, \textit{prompt} & BPMN XML & PMo, 46 & 97,8\% (XSD) & --- & juiz 3,77 \\
			Proposta, A2r & DeepSeek~V4~Pro, \textit{prompt} & BPMN-DSL & PMo, 46 & 87,0\% (XSD) & 84,0\% & juiz 3,85 \\
			Proposta, A4 & Qwen2.5-Coder-7B, QLoRA & BPMN-DSL & PMo, 46 & 87,0\% (XSD) & 82,7\% & juiz 2,78 \\
			\bottomrule
		\end{tabular}%
	}{%
	\Fonte{Elaborado pelo autor, com os valores reportados em cada trabalho}%
	\Nota{Validade: proporção de saídas válidas no critério indicado entre parênteses. XSD: aceita pelo esquema do BPMN 2.0; desenhável: sem erro de sintaxe nem fluxo para elemento inexistente; sintaxe: lida pelo analisador da linguagem de saída. Redução de \textit{tokens}: proporção de \textit{tokens} a menos que o BPMN XML, medida nos modelos de referência em \citeonline{brissard2025pmr}, na saída média dos modelos em \citeonline{licardo2026bpmn} e pela razão de compressão na proposta. Qualidade: medida própria de cada trabalho, e os valores não são comparáveis entre linhas. Na proposta, nota de fidelidade do juiz, de 0 a 5, exploratória (Seção~\ref{subsec:fidelidade}). \textsuperscript{a}~Em relação ao BPMN XML com as coordenadas do diagrama; sem elas, a redução é de 70\%. \textsuperscript{b}~Modelo ajustado, com teste no mesmo conjunto do treino. \textsuperscript{c}~Em linhas de código, e não em \textit{tokens}. \textsuperscript{d}~O artigo não deixa claro se os 88\% se referem à geração direta de BPMN ou ao BPMN obtido a partir do DOT.}%
}
\end{table}

A comparação direta só é possível em duas colunas. A validade pelo esquema XSD é medida da mesma forma em \citeonline{lauer2026assessing} e na proposta, e nos dois casos a geração direta de XML por modelos de grande porte fica acima de 97\%. A redução de \textit{tokens} é medida sobre o mesmo conjunto em \citeonline{brissard2025pmr} e na proposta, com duas ressalvas: a referência deles é o XML com as coordenadas do diagrama, que a proposta exclui, e os tokenizadores são diferentes. Tomando como base o XML sem coordenadas, que é a base da proposta, o Mermaid, a representação mais compacta do estudo, tem cerca de 76\% menos \textit{tokens}, e a BPMN-DSL, entre 81\% e 84\%, conforme o conjunto. As demais medidas de validade usam critérios mais fracos que o esquema, e as medidas de qualidade são todas diferentes, cada uma com seu próprio conjunto de teste, de modo que os números da última coluna não podem ser ordenados entre trabalhos.

Três aspectos distinguem a proposta dos trabalhos analisados. O primeiro é o custo: os trabalhos medem \textit{tokens} ou latência, mas nenhum converte essas medidas em custo por requisição nem desconta as saídas inválidas, como faz o custo por saída válida (Seção~\ref{sec:res-custo}). O segundo é a combinação entre representação compacta e modelo pequeno ajustado: os trabalhos com representação intermediária usam modelos de grande porte por \textit{prompt}, e os que ajustam modelos pequenos geram DOT. Entre estes, \citeonline{li2025llmbased} e \citeonline{celikmasat2026instrubpm} avaliam a saída na própria linguagem DOT, sem medir a compressão em relação ao XML nem validar pelo esquema o BPMN resultante. O terceiro é a condição de teste: nos trabalhos com modelos ajustados, o teste vem do mesmo conjunto que o treino, enquanto na proposta o treino e o teste vêm de fontes diferentes, e a fidelidade é julgada contra o texto, e não contra um modelo de referência, por um juiz cuja validação é reportada junto com os resultados.
